\section*{अध्याय \ref{ch:b3:total-synthesis} --- पूर्ण संश्लेषण: रणनीति और क्लासिकी मार्ग}

\begin{solution}{exo:b3:total-synthesis:1}
90\,\% पर: $0.9^{10} = 35$\,\% और $0.9^{20} = 12$\,\%। 80\,\% पर: $0.8^{10} = 11$\,\% और $0.8^{20} = 1.2$\,\%।
\end{solution}

\begin{solution}{exo:b3:total-synthesis:2}
कुल $5 + 3 + 1 + 4 = 13$ चरण; LLS $5 + 1 + 4 = 10$।
\end{solution}

\begin{solution}{exo:b3:total-synthesis:3}
4-(डाइमेथिलऐमीनो)ब्यूटेन-2-ओन, $\ce{CH3COCH2CH2N(CH3)2}$: ऐसीटोन का ईनॉल इमीनियम
आयन $\ce{CH2=N+(CH3)2}$ पर जुड़ता है।
\end{solution}

\begin{solution}{exo:b3:total-synthesis:4}
C=C के ऐलिलिक दोनों वलय-आबंधों का विच्छेदन कीजिए: ब्यूटा-1,3-डाईन और ब्यूट-3-ईन-2-ओन
(मेथिल वाइनिल कीटोन), जो गर्म करने पर अभिक्रिया करते हैं।
\end{solution}

\begin{solution}{exo:b3:total-synthesis:5}
रैखिक: $0.9^{21} = 11$\,\%। अभिसारी: $0.9^{11} = 31$\,\%, लगभग तीन गुना।
\end{solution}

\begin{solution}{exo:b3:total-synthesis:6}
A: $(7 + 1)/12 = 67$\,\%। B: $7/9 = 78$\,\%।
\end{solution}

\begin{solution}{exo:b3:total-synthesis:7}
उनके बिना: $0.9^8 = 43$\,\%। 95\,\% वाले चार अतिरिक्त चरणों के साथ: $43\,\% \times 0.95^4 = 35$\,\%।
\end{solution}

\begin{solution}{exo:b3:total-synthesis:8}\emergencystretch=3em
प्राथमिक: एक भारी सिलिल ईथर (\textit{tert}-ब्यूटिलडाइफ़ेनिलसिलिल), फ़्लुओराइड से हटाया जाता है।
द्वितीयक: एक 4-मेथॉक्सीबेंज़िल ईथर, DDQ से ऑक्सीकरण द्वारा हटाया जाता है। बाधित
तृतीयक ऐल्कोहॉल को प्रायः मुक्त छोड़ा जा सकता है, या एक छोटे सिलिल ईथर के रूप में रक्षित
किया जा सकता है जो मृदु अम्ल से विदलित होता है। प्रत्येक परिस्थिति अन्य समूहों को अक्षत छोड़ती है।
\end{solution}

\begin{solution}{exo:b3:total-synthesis:9}
\omterm{def:b3:total-synthesis:total}{पूर्ण संश्लेषणों} में दसियों चरण लगते थे और \omterm{def:b3:total-synthesis:architecture}{समग्र लब्धियाँ} 1\,\% से बहुत नीचे थीं; पूर्ण
वलय-तंत्र और अधिकांश त्रिविमजनक केंद्रों वाला एक यौगिक, जो एक नवीकरणीय स्रोत, यूरोपीय यू
की पत्तियों से निष्कर्षित होता है, कुछ ही चरणों में रूपांतरित हो जाता है।
\end{solution}

\begin{solution}{exo:b3:total-synthesis:10}
माइकल योग: डाइओन का ईनोलेट (C2, दोनों कार्बोनिलों के बीच) मेथिल वाइनिल कीटोन के
$\beta$-कार्बन पर जुड़ता है, जिससे एक ट्राइकीटोन बनता है। ऐल्डोल: मेथिल कीटोन का
ईनोलेट अंतःआण्विक रूप से एक वलय-कार्बोनिल पर आक्रमण करता है, और निर्जलीकरण
साइक्लोहेक्सीनोन देता है: वीलैंड--मीशर कीटोन, कोणीय मेथिल समूह वाला
एक द्विचक्रीय ईनडाइओन।
\end{solution}

\begin{solution}{exo:b3:total-synthesis:11}
$y^{m+1}/y^{2m+1} = y^{-m}$: 90\,\% पर, $m = 5$ के लिए 1.7 और $m = 15$ के लिए 4.9। मार्ग जितना लंबा, अभिसरण
उतना ही लाभदायक।
\end{solution}

\begin{solution}{exo:b3:total-synthesis:12}
प्रत्येक वलय धनायन और आक्रमणकारी द्विबंध की कुर्सी-सदृश व्यवस्था में
बंद होता है; प्रत्येक द्विबंध पर योग anti होता है, और प्रतिस्थापी
छद्म-विषुवतीय स्थितियाँ लेते हैं, जिससे प्रत्येक वलय-संगलन \textit{trans} होता है। \textit{E} द्विबंध
अपनी दोनों शृंखला-निरंतरताओं को ऐसे रखता है कि \textit{trans} संगलन बने; \textit{Z} द्विबंध
\textit{cis} संगलन देता: पॉलिईन की ज्यामिति उत्पाद के त्रिविम रसायन को कूटबद्ध करती है।
\end{solution}

\begin{solution}{pb:b3:total-synthesis:1}
\textbf{1.} $\ce{C4H6O2}$, $\ce{CH5N}$, $\ce{C5H6O5}$, $\ce{C8H13NO}$।

\textbf{2.} $\ce{C4H6O2 + CH5N + C5H6O5 -> C8H13NO + 2CO2 + 2H2O}$।

\textbf{3.} 86.0, 31.0 और \qty{146.0}{g/mol}; ट्रोपिनोन \qty{139.0}{g/mol}।

\textbf{4.} $139.0/(86.0 + 31.0 + 146.0) = 0.53$: 53\,\%।

\textbf{5.} इसके केंद्रीय $\ce{CH2}$ समूह, जिनमें से प्रत्येक दो कार्बोनिल समूहों के बीच है, जल में
सरलता से ईनॉलित होते हैं; ऐसीटोन स्वयं बहुत दुर्बल नाभिकरागी है, और सक्रियक कार्बोक्सिल
समूह बाद में निकल जाते हैं।

\textbf{6.} इमीनियम आयनों को मुक्त ऐमीन चाहिए (अधिक अम्ल उसे प्रोटॉनित कर देता है), और
ईनॉल तथा इमीनियम के बनने को कुछ अम्ल चाहिए: उदासीनता के निकट pH इन्हें संतुलित करता है।

\textbf{7.} $\ce{CH3NH2}$ एक ऐल्डिहाइड समूह पर जुड़ता है; जल का निकास इमीनियम आयन
$\ce{R-CH=N+H-CH3}$ देता है।

\textbf{8.} ऐसीटोनडाइकार्बोक्सिलिक अम्ल का एक ईनॉल कार्बन इमीनियम कार्बन पर आक्रमण करता है, जिससे
एक C--C आबंध और एक द्वितीयक ऐमीन बनते हैं।

\textbf{9.} द्वितीयक ऐमीन उसी शृंखला के दूसरे ऐल्डिहाइड समूह के साथ संघनित होता है,
जिससे एक चक्रीय (पंचसदस्यीय) इमीनियम आयन बनता है।

\textbf{10.} कीटोन के दूसरी ओर का ईनॉल उस पर आक्रमण करता है, जिससे
छह-सदस्यीय वलय और द्विचक्रीय ढाँचा बंद होता है।

\textbf{11.} प्रत्येक एक $\beta$-कीटो अम्ल है, जो एक चक्रीय छह-सदस्यीय संक्रमण
अवस्था से होकर $\ce{CO2}$ खोता है, जिससे कीटोन का ईनॉल मिलता है।

\textbf{12.} दो C--C आबंध और दो C--N आबंध।

\textbf{13.} $0.75^{14} = 0.018$: 1.8\,\%।

\textbf{14.} $42/1.78 = 24$।

\textbf{15.} 100\,\%: एक चरण, सब कुछ आबंध-निर्माण।

\textbf{16.} एक संक्रिया, विलायक के रूप में जल, सस्ते घटक, कोई रक्षी समूह नहीं और
कार्बन डाइऑक्साइड तथा जल के अतिरिक्त लगभग कोई अपशिष्ट नहीं।

\textbf{17.} पौधा ट्रोपेन ढाँचे को एक ऐमीनो-अम्ल-व्युत्पन्न
चक्रीय इमीनियम आयन और एक ऐसीटेट-व्युत्पन्न इकाई से, उसी प्रकार के मैनिख
रसायन द्वारा, बनाता है।

\textbf{18.} \omterm{def:b3:total-synthesis:cascade}{सोपानी अभिक्रियाएँ} (और डील्स--ऐल्डर अभिक्रिया जैसे वलय बनाने वाले
\omterm{def:b3:pericyclic:families}{चक्रीयोजन})।

\textbf{19.} नहीं: एक \omterm{def:b3:point-groups:mirror}{दर्पण तल} N, उसके मेथिल समूह, C3 और कार्बोनिल
ऑक्सीजन से होकर जाता है।

\textbf{20.} प्रत्येक पर चार भिन्न समूह हैं, पर वे अणु के अपने सममिति तल में
एक-दूसरे के दर्पण प्रतिबिंब हैं: एक मेसो-प्रकार का यौगिक।

\textbf{21.} अप्रतिबिंबरूपी, जो नाइट्रोजन सेतु के सापेक्ष C3 पर OH के
अभिविन्यास में भिन्न हैं।

\textbf{22.} नहीं: दोनों में \omterm{def:b3:point-groups:mirror}{दर्पण तल} बना रहता है।

\textbf{23.} कार्बोनिल समूह के दोनों फलक तुल्य नहीं हैं: एक नाइट्रोजन सेतु
की ओर है, दूसरा दो-कार्बन सेतु की ओर, और ये भिन्न रूप से बाधा डालते हैं।

\textbf{24.} एक-पात्र मार्ग 14-चरणीय रैखिक मार्ग की \omterm{def:b3:total-synthesis:architecture}{समग्र लब्धि} का लगभग 24 गुना
देता है।
\end{solution}
